\capitulo{CAPÍTULO 8}{La reparación directa como vía procesal para reclamar la responsabilidad extracontractual del Estado}

\section{8.1. El medio de control de reparación directa (art. 140 de la Ley 1437 de 2011)}
\textbf{Finalidad:} la reparación del daño antijurídico.

Admite tres esquemas de partes:

\begin{tabla}
\begin{longtable}{|>{\raggedright\arraybackslash}m{\dimexpr 0.245\linewidth-2\tabcolsep-0.600pt\relax}|>{\raggedright\arraybackslash}m{\dimexpr 0.755\linewidth-2\tabcolsep-0.600pt\relax}|}
\hline
\cellcolor[HTML]{1F3864}\textcolor[HTML]{FFFFFF}{\textbf{Demandante}} & \cellcolor[HTML]{1F3864}\textcolor[HTML]{FFFFFF}{\textbf{Demandado}} \\ \hline
Persona interesada & Entidad pública causante del daño \\ \hline
\cellcolor[HTML]{F2F5FA}Entidad pública perjudicada & \cellcolor[HTML]{F2F5FA}Particular / entidad causante del daño \\ \hline
Persona interesada & Particular que haya obrado siguiendo una expresa instrucción de una entidad pública \\ \hline
\end{longtable}
\end{tabla}
\begin{comentario}
\textcolor[HTML]{8A6D3B}{\textbf{Comentario. }}Sobre el alcance de la expresión «particular que haya obrado siguiendo una expresa instrucción» y su interpretación en la sentencia C-644 de 2011, así como sobre el último inciso del art. 140 (proporción de responsabilidad cuando concurren particulares y entidades públicas), ver Capítulo 5, secciones 5.10 y 5.11.4.
\end{comentario}
\begin{ampliacion}
\textcolor[HTML]{5B3F8C}{\textbf{Ampliación. }}El art. 140 del CPACA regula el medio de control de reparación directa: «la persona interesada podrá demandar directamente la reparación del daño antijurídico producido por la acción u omisión de los agentes del Estado». Es el cauce procesal del art. 90 de la Constitución para los daños que \textbf{no} provienen de un acto administrativo (cuando el daño lo causa un acto, hay que demandarlo por nulidad y restablecimiento del derecho). La norma también permite que las entidades públicas demanden a particulares o a otras entidades que les causen daños, y, según su último inciso, cuando concurren particulares y entidades públicas en la causación del daño, la sentencia determina la proporción por la que responde cada uno según su influencia causal.
\end{ampliacion}
\section{8.2. Fuente del daño}
\begin{itemize}
\item Hecho administrativo.
\item Omisión administrativa.
\item Operación administrativa.
\item Ocupación temporal o permanente de inmuebles por obras públicas.
\item Otras causas imputables a entidades públicas o a particulares que obren siguiendo instrucciones de entidades públicas.
\end{itemize}
\section{8.3. ¿Qué se puede pedir?}
\begin{itemize}
\item Reparación de daños antijurídicos de origen extracontractual.
\item Responsabilidad precontractual.
\item \textit{Actio in rem verso}.
\item Daño al patrimonio público por actos de corrupción.
\end{itemize}
\section{8.4. Legitimación en la causa por activa}
\begin{itemize}
\item El particular con interés.
\item Las entidades públicas perjudicadas.
\end{itemize}
\begin{comentario}
\textcolor[HTML]{8A6D3B}{\textbf{Comentario. }}Sobre la legitimación de las víctimas directas e indirectas, la reclamación \textit{iure proprio} o \textit{iure hereditario} y la legitimación del propietario, poseedor o tenedor en los daños a las cosas, ver Capítulo 4, secciones 4.4 y 4.5.3.
\end{comentario}
\section{8.5. Caducidad}
\begin{itemize}
\item Dos (2) años contados a partir del día siguiente al de la ocurrencia de la acción u omisión causante del daño.
\item Dos (2) años contados a partir de cuando el demandante tuvo o debió tener conocimiento del mismo, si fue en fecha posterior y siempre que pruebe la imposibilidad de haberlo conocido en la fecha de su ocurrencia.
\item Delito de desaparición forzada: dos (2) años contados a partir de la fecha en que aparezca la víctima o, en su defecto, desde la ejecutoria del fallo definitivo adoptado en el proceso penal, sin perjuicio de que la demanda con tal pretensión pueda intentarse desde el momento en que ocurrieron los hechos que dieron lugar a la desaparición.
\item Daños causados por la comisión de delitos de lesa humanidad: sentencia de unificación del 29 de enero de 2020, expediente 61033.
\end{itemize}
\begin{comentario}
\textcolor[HTML]{8A6D3B}{\textbf{Comentario. }}El desarrollo de estas reglas (el conocimiento del daño y la carga de invocarlo desde la demanda, los casos prácticos, la caducidad en lesa humanidad y el debate sobre el desplazamiento forzado) está en el Capítulo 4, sección 4.3. Las reglas de caducidad en el error judicial cuando se interpuso tutela están en el Capítulo 7, sección 7.2.3.
\end{comentario}
\begin{ampliacion}
\textcolor[HTML]{5B3F8C}{\textbf{Ampliación. }}Esta regla está en el art. 164, num. 2, lit. i) del CPACA. La caducidad es un presupuesto procesal de orden público: el término no se suspende ni se interrumpe por voluntad de las partes, salvo la suspensión que produce la solicitud de conciliación extrajudicial, y el juez debe declararla incluso de oficio (por eso, si es evidente, rechaza la demanda). La conciliación extrajudicial, cuando el asunto es conciliable, es requisito de procedibilidad de la reparación directa (art. 161, num. 1, del CPACA); la Ley 2195 de 2022 exceptúa de ese requisito a la entidad que demanda la reparación del daño causado por actos de corrupción (ver Capítulo 9).
\end{ampliacion}
\section{8.6. Otras reglas procesales relevantes}
\begin{itemize}
\item \textbf{Fuero de atracción:} cuando el daño es causado por el Estado y por un particular, se puede demandar a ambos ante la jurisdicción contenciosa, que conoce de la responsabilidad de los dos (Capítulo 5, sección 5.6.1).
\item \textbf{Estimación razonable de la cuantía:} es requisito de la demanda (art. 162 del CPACA); es el primer momento en que se cuantifican los perjuicios (Capítulo 4, sección 4.7.2).
\item \textbf{Condena en abstracto e incidente de liquidación:} cuando se prueba el daño pero no la cuantía del perjuicio (Capítulo 4, sección 4.5.1).
\item \textbf{Prueba sobreviniente:} por ejemplo, el dictamen de la Junta Médico Laboral expedido durante el proceso (Capítulo 4, sección 4.5.1).
\item \textbf{Costas:} si se pierde el proceso, la condena en costas es proporcional a la cuantía de los perjuicios reclamados (Capítulo 4, sección 4.9.6).
\item \textbf{Dictamen pericial:} necesario, por regla general, en responsabilidad médica (Capítulo 5, sección 5.6.2).
\end{itemize}
